\documentclass[11pt]{article}

\usepackage[margin=1in]{geometry}
\usepackage[T1]{fontenc}
\usepackage{amsmath, amssymb}
\usepackage{booktabs}
\usepackage{enumitem}
\usepackage{titlesec}
\usepackage{fancyhdr}
\usepackage{float}
\usepackage{xcolor}

\usepackage[hidelinks]{hyperref}

\titlespacing*{\section}{0pt}{1.2em}{0.5em}
\titlespacing*{\subsection}{0pt}{0.8em}{0.4em}
\titleformat{\section}{\large\bfseries}{\thesection.}{0.5em}{}
\titleformat{\subsection}{\normalsize\bfseries}{\thesubsection}{0.5em}{}

\pagestyle{fancy}
\fancyhf{}
\lhead{\small CS 388: Natural Language Processing}
\rhead{\small Homework 1}
\cfoot{\small\thepage}
\renewcommand{\headrulewidth}{0.4pt}

\setlist[enumerate]{itemsep=2pt, topsep=3pt}
\setlist[itemize]{itemsep=2pt, topsep=3pt}

\newcommand{\bs}{\ensuremath{\langle \textsc{s} \rangle}}
\newcommand{\es}{\ensuremath{\langle \textsc{/s} \rangle}}
\newcommand{\Count}[1]{\ensuremath{C(#1)}}
\newcommand{\wrd}[1]{\texttt{#1}}

\definecolor{solblue}{RGB}{0,60,180}
\newcommand{\sol}[1]{\par\medskip{\color{solblue}\noindent\textbf{Solution.} #1}\par\medskip}
\newcommand{\B}[1]{{\color{solblue}#1}}

\begin{document}

\begin{center}
  {\Large\bfseries Homework 1: $n$-gram Language Models}\\[2pt]
  {\normalsize CS 388: Natural Language Processing \quad\textbullet\quad Fall 2026}\\[2pt]
\end{center}
\vspace{0.5em}
\hrule
\vspace{1em}

\noindent
\textbf{Instructions.} This assignment can be done entirely by hand -- no coding required; however, please transcribe your answers/work into the provided \LaTeX file and submit the PDF, rather than a handwritten response.
\textbf{Please use the {\tt{\textbackslash B\{\}}} command to wrap your solutions, making them blue (this will help us when grading).}
Show your work. Report probabilities as exact fractions (e.g.\ $\tfrac{3}{14}$)
rather than decimals, and simplify where convenient. Submit a single PDF to Gradescope.

\section{Setup and Notation}

Throughout, we use the following corpus of $10$ sentences as our training data. Assume the text has already
been tokenized and lowercased, so each whitespace-separated string is one token.

\begin{center}
\begin{tabular}{rl@{\hskip 3em}rl}
\toprule
\multicolumn{4}{c}{\textbf{Training corpus} $\mathcal{D}$}\\
\midrule
1. & \wrd{the cat chased the dog}   & 6.  & \wrd{a dog saw the cat}\\
2. & \wrd{the dog chased the cat}   & 7.  & \wrd{the mouse saw the cat}\\
3. & \wrd{the cat sat on the mat}   & 8.  & \wrd{the mouse slept}\\
4. & \wrd{the dog sat on the mat}   & 9.  & \wrd{the cat slept}\\
5. & \wrd{a cat saw the dog}        & 10. & \wrd{a dog barked}\\
\bottomrule
\end{tabular}
\end{center}

\noindent
\textbf{Padding.} Every sentence is padded with sentence-boundary symbols. For the \emph{bigram} model, pad each
sentence with one \bs{} on the left and one \es{} on the right; for the \emph{trigram} model, pad with two \bs{}
on the left and one \es{} on the right. So sentence 8 is treated as
\[
\bs\ \bs\ \wrd{the}\ \wrd{mouse}\ \wrd{slept}\ \es .
\]
The symbol \es{} is predicted like any other token (this is what makes the model a proper distribution over
sentences), but \bs{} is never predicted -- it only ever appears in a conditioning context.

\medskip\noindent
\textbf{Vocabulary.} Let $V$ denote the vocabulary of \emph{possible next tokens}: the $12$ word types that occur in
$\mathcal{D}$, plus \es{}, and \emph{not} \bs{}. Hence $|V| = 13$. Use this value wherever $|V|$ appears.



\medskip\noindent
\textbf{Convention.} We consider any probability of the form $\frac{n}{0}$ to be undefined.


\section{Counting (30 points)}
In class we saw the maximum likelihood estimate (MLE) for n-gram models, based on counting occurrences. 
Your first task is to build the count tables you will need for this estimator. Include the padding symbols in your counts.

\begin{enumerate}[label=(\alph*)]

  \item \textbf{(25 pts)} Report the full bigram and trigram count tables. Tuples with count zero may be omitted.
  
 Example (bigram)

 \begin{table}[H]
     \centering
     \begin{tabular}{c|c}
     \toprule
     Bigram & Count \\
     \midrule
     \wrd{the}\, \wrd{cat} & 6\\
      ... & ... \\
      \bottomrule
     \end{tabular}
 \end{table}

  \item \textbf{(5 pts)} Why do we need the padding tokens \bs{} and \es{}? What happens if we do not
  include them?

\end{enumerate}

\section{Sentence Probabilities (40 points)}

Consider the following three test sentences.

\begin{center}
\begin{tabular}{ll}
\toprule
$S_1$ & \wrd{the dog chased the dog}\\
$S_2$ & \wrd{the cat chased the mouse}\\
$S_3$ & \wrd{the mouse chased the dog}\\
\bottomrule
\end{tabular}
\end{center}

For each part below, write out the full product of conditional probabilities, and their corresponding values, as well as the final product. 

\begin{enumerate}[label=(\alph*)]
  \item \textbf{(12 pts)} Compute $P_{\text{bi}}(S_1)$ and $P_{\text{tri}}(S_1)$.

  \item \textbf{(14 pts)} Compute $P_{\text{bi}}(S_2)$ and $P_{\text{tri}}(S_2)$.

  \item \textbf{(14 pts)} Compute $P_{\text{bi}}(S_3)$ and $P_{\text{tri}}(S_3)$.

\end{enumerate}

\section{Laplace Smoothing and Backoff (30 points)}

In class we saw both Laplace smoothing and Backoff. As a reminder:

\medskip\noindent
\textbf{Add-$k$ (Laplace) smoothing, $k=1$.} Add one pseudo-count to every possible continuation:
\[
P^{+1}_{\text{bi}}(w_i \mid w_{i-1}) = \frac{\Count{w_{i-1} w_i} + 1}{\Count{w_{i-1}} + V},
\qquad
P^{+1}_{\text{tri}}(w_i \mid w_{i-2} w_{i-1}) = \frac{\Count{w_{i-2} w_{i-1} w_i} + 1}{\Count{w_{i-2} w_{i-1}} + V}.
\]

\medskip\noindent
\textbf{Backoff.} Adding $1$ to the numerator does not help address problems stemming from contexts we have never observed. This is especially problematic as we increase $n$: for a trigram model, there will be many contexts we have never seen. 
In this case, we can back-off an $n$-gram to an $n-1$-gram, if we have not seen the particular $n$-gram in question. We can repeat this until we get to unigrams.

For this assignment, to compute backoff , first use the Laplace-smoothed trigram estimate whenever the trigram conditioning context has been observed. If the bigram context $w_{i-2}w_{i-1}$ has never been observed, i.e.\ $\Count{w_{i-2}w_{i-1}} = 0$,
we back off to the Laplace-smoothed bigram estimate. \textbf{Smoothing is therefore
applied in all cases}:
\[
P_{\text{bo}}(w_i \mid w_{i-2} w_{i-1}) =
\begin{cases}
\dfrac{\Count{w_{i-2} w_{i-1} w_i} + 1}{\Count{w_{i-2} w_{i-1}} + V}, & \text{if } \Count{w_{i-2} w_{i-1}} > 0, \\[12pt]
\dfrac{\Count{w_{i-1} w_i} + 1}{\Count{w_{i-1}} + V},                 & \text{if } \Count{w_{i-2} w_{i-1}} = 0.
\end{cases}
\]


\begin{enumerate}[label=(\alph*)]
  \item \textbf{(11 pts)} Recompute $P^{+1}_{\text{bi}}(S_2)$ and $P_{\text{bo}}(S_2)$, again showing every
  factor.

  \item \textbf{(11 pts)} Recompute $P^{+1}_{\text{bi}}(S_3)$ and $P_{\text{bo}}(S_3)$.


  \item \textbf{(8 pts)} Recompute $P^{+1}_{\text{bi}}(S_1)$ and compare it to $P_{\text{bi}}(S_1)$ from Part 2. What difference do you observe, and what is the explanation for the direction of that difference? 

\end{enumerate}

\vspace{1em}
\hrule
\vspace{0.5em}
\noindent
\small
\textbf{What to turn in.} A single PDF containing your count tables (Part 1), the full factor-by-factor
derivations for all six probability computations in Part 2 and all five in Part 3, and your written answers to
the discussion questions. 

\end{document}
